% ID: SEL-D0013-BD-E07
% ESTADO: cotejada. Segunda pasada visual completa y borrador compilado revisados.
% Fuente: Selectividad/matematicas_aplicadas/matematicas_aplicadas_P_A.pdf
% SHA-256: 0305750f6b6c3cc3edbba71a19cd7cce4276935dff5a0b7d964b02eaedeaa222
% Página: 2; bloque D; ejercicio original 7.

\begin{enumerate}[label=\alph*)]
  \item \textbf{(1 punto)} Una población de $25\,000$ personas se ha dividido
  en cuatro estratos con tamaños $15\,000$, $5\,000$, $3\,000$ y $2\,000$
  personas respectivamente. En esa población se ha realizado un muestreo
  estratificado con afijación proporcional, en el que se han elegido al azar
  36 personas del tercer estrato. Determine el tamaño de la muestra total
  obtenida con este muestreo y su composición.
  \item \textbf{(1.5 puntos)} Dada la población $P=\{2,4,6\}$, construya todas
  las muestras posibles de tamaño 2 que se puedan formar mediante muestreo
  aleatorio simple y halle la desviación típica de las medias muestrales
  obtenidas con todas esas muestras.
\end{enumerate}
